% Generated from THESIS_DRAFT.md.
\chapter*{Tóm tắt}
\thispagestyle{empty}

Trong thao tác robot ở chuỗi tác vụ dài, các quan sát cảm biến robot thu được sau khi hoàn tất một chu kỳ thao tác thường tương tự với quan sát ban đầu, trong khi quyết định tiếp tục hay dừng chuỗi hành động lại phụ thuộc vào số chu kỳ và mốc tiến độ đã hoàn tất. Nhằm giải quyết hiện tượng này, khóa luận đề xuất một kiến trúc bộ nhớ theo sự kiện có khả năng biểu diễn mạch lạc thông tin tiến độ và hỗ trợ các mô hình thị giác - ngôn ngữ - hành động (VLA) thực hiện các tác vụ chuỗi cần bộ nhớ. Nghiên cứu tập trung giải quyết bài toán trên bộ tiêu chuẩn thao tác chuỗi dài, đặc biệt hướng tới việc nâng cao năng lực cho các mô hình VLA nhỏ gọn vốn chịu nhiều ràng buộc nghiêm ngặt về độ dài ngữ cảnh và tài nguyên tính toán.

Phương pháp đề xuất biểu diễn tiến độ bằng cách xác định trạng thái vật thể và robot hiện tại, mô tả lại sự kiện đã xảy ra một cách thích ứng với tác vụ và quan sát. Mỗi sự kiện đều được lưu vết đầy đủ bằng chứng và thời điểm để sử dụng nhằm đảm bảo độ tin cậy và tính nhân quả. Ngay cả khi lịch sử chưa được quan sát đầy đủ, bộ nhớ vẫn có thể biểu diễn tiến độ một cách thích ứng. Hai tính chất của hệ bộ nhớ đề xuất được chứng minh trong phạm vi khóa luận bao gồm: Bộ đếm và ghi nhận sự kiện chỉ thay đổi khi xuất hiện một sự kiện kích hoạt hợp lệ; Nội dung cấp cho chính sách tại một thời điểm chỉ phụ thuộc vào dữ liệu sẵn có đến thời điểm đó. Bộ nhớ sự kiện này sau đó được chiếu thành một bản tóm tắt ngôn ngữ tự nhiên gọn nhẹ, đưa vào đầu vào điều kiện hóa của mô hình VLA với phần sinh hành động đã được huấn luyện thích ứng.

Hiệu quả của phương pháp được kiểm chứng thông qua các thực nghiệm trên Bộ tiêu chuẩn đánh giá khả năng học tập suốt đời của robot LIBERO và Bộ tiêu chuẩn thao tác chuỗi dài LIBERO-Mem, bao gồm cả đánh giá về năng lực tri giác và đánh giá điều khiển vòng kín trên mô hình VLA. Kết quả cho thấy bộ theo dõi thị giác có khả năng phát hiện và đồng bộ chính xác các mốc sự kiện vật lý từ dòng quan sát thô mà không cần tiếp cận trạng thái đặc quyền của môi trường mô phỏng. Khi tích hợp vào vòng điều khiển kín, thông tin bộ nhớ sự kiện cung cấp tín hiệu tiến độ then chốt giúp mô hình phân giải tính đa nghĩa của quan sát thị giác, định hướng chuỗi hành động đúng tiến độ và ngăn ngừa các lỗi lặp lại chu kỳ hoặc kết thúc sớm. Đồng thời, các phân tích chuyên sâu làm sáng tỏ mối quan hệ phụ thuộc giữa độ trễ tích lũy bằng chứng, độ hoàn thiện của lịch sử quan sát và độ ổn định của chính sách điều khiển. Nghiên cứu qua đó xác lập một khung phương pháp luận có tính khả kiểm cao, đặt nền tảng vững chắc cho việc mở rộng kiến trúc bộ nhớ sự kiện - bằng chứng cho các bài toán thao tác đa bước phức tạp trên các mô hình robot học sâu.

\textbf{Từ khóa:} mô hình thị giác - ngôn ngữ - hành động, SmolVLA, bộ nhớ sự kiện, quan sát một phần, thao tác robot.

